% !TeX program = xelatex
\documentclass[aspectratio=169]{beamer}
\usetheme{thubeamer}

\usepackage{booktabs}
\usepackage{amsmath}
\usepackage{bm}

\graphicspath{{figures/}}
% Local macros for thubeamer slides
\newcommand{\vx}{\mathbf{x}}
\newcommand{\reals}{\mathbb{R}}


\newcommand{\T}{\mathrm{T}}
\newcommand{\R}{\mathbb{R}}

\title[地下停车场应急人员定位]{地下停车场应急人员定位问题研究\\[2mm]基于蓝牙 RSSI 与步行约束融合的建模与求解}
\author[陆雨澍]{学生：陆雨澍\\[5mm] 导师：陈务深~教授}
\institute[南京工业大学]{\small 南京工业大学\\数理科学学院数学系}
\date{\small \vskip -10pt 2026 年 4 月}

\begin{document}

\begin{frame}
  \maketitle
\end{frame}

\section*{目录}
\frame{
  \frametitle{\secname}
  \tableofcontents[hideallsubsections]
}

\section{研究背景}

\subsection{问题来源}

\begin{frame}{问题来源}
  \begin{block}{应用场景}
    \begin{itemize}
      \setlength{\itemsep}{6pt}
      \item 地下停车场火灾、烟雾扩散或结构异常时，GNSS 信号不可用
      \item 人工语音定位存在延迟与误判
      \item BLE 信标部署成本低，适合规则空间应急定位
    \end{itemize}
  \end{block}
  \begin{block}{研究目标}
    在噪声与缺失观测下，估计人员轨迹 $\{\bm p_t\}_{t=1}^{T}$，$\bm p_t=[x_t,y_t]^{\T}$
  \end{block}
\end{frame}

\begin{frame}{研究难点}
  \begin{enumerate}
    \setlength{\itemsep}{8pt}
    \item 金属车体与立柱导致多路径反射，RSSI 波动大
    \item 车辆动态进出造成遮挡，观测噪声非平稳
    \item 单时刻点位估计抖动明显，难以形成可用轨迹
    \item 应急场景要求算法实时、稳定、可解释
  \end{enumerate}
\end{frame}

\begin{frame}{技术路线}
  \begin{enumerate}
    \setlength{\itemsep}{8pt}
    \item 对数路径损耗模型：RSSI $\rightarrow$ 距离观测
    \item 多信标加权最小二乘：时刻级粗定位
    \item 步长与转向约束：滑动窗口轨迹平滑
    \item 迭代重加权（IRWLS）：抑制粗差，输出稳定轨迹
  \end{enumerate}
  \vspace{6mm}
  \centering
  场景抽象 $\rightarrow$ 数学建模 $\rightarrow$ 鲁棒优化 $\rightarrow$ 实验评估
\end{frame}

\section{数学建模}

\subsection{观测模型}

\begin{frame}{RSSI 路径损耗模型}
  \begin{block}{观测方程}
    \[
      z_{i,t}=P_0-10n\log_{10}(d_{i,t})+\varepsilon_{i,t}
    \]
  \end{block}
  \begin{block}{距离反演}
    \[
      \hat d_{i,t}=10^{\frac{P_0-z_{i,t}}{10n}}
    \]
  \end{block}
  \begin{itemize}
    \item $P_0$：参考距离处信号强度；$n$：路径损耗指数
    \item $d_{i,t}$：人员到第 $i$ 个信标的几何距离
  \end{itemize}
\end{frame}

\begin{frame}{几何观测关系}
  信标坐标 $\bm b_i=[u_i,v_i]^{\T}$，理论距离
  \[
    d_{i,t}(\bm p_t)=\sqrt{(x_t-u_i)^2+(y_t-v_i)^2}
  \]
  观测残差
  \[
    r_{i,t}=d_{i,t}(\bm p_t)-\hat d_{i,t}
  \]
  \begin{block}{单帧目标函数}
    \[
      F(\bm p)=\sum_i \omega_i r_i(\bm p)^2
    \]
  \end{block}
\end{frame}

\begin{frame}{融合目标函数}
  滑动窗口 $\mathcal W_t=\{t-w+1,\dots,t\}$ 上：
  \small
  \begin{align*}
    J=&\sum_{\tau\in\mathcal W_t}\sum_{i\in\mathcal B_\tau}
    \omega_{i,\tau}\left(d_{i,\tau}(\bm p_\tau)-\hat d_{i,\tau}\right)^2 \\
    &+\lambda_s\sum_{\tau\in\mathcal W_t}\|\bm p_\tau-\bm p_{\tau-1}\|_2^2
    +\lambda_h\sum_{\tau\in\mathcal W_t}\phi(\Delta\theta_\tau)
  \end{align*}
  \normalsize
  \begin{itemize}
    \item $\omega_{i,\tau}$：Huber 鲁棒权重
    \item $\lambda_s$：步长平滑；$\lambda_h$：航向平滑
  \end{itemize}
\end{frame}

\begin{frame}{鲁棒权重与可行域约束}
  \begin{columns}[T]
    \column{0.48\textwidth}
    \begin{block}{Huber 权重}
      \[
        \omega_{i,\tau}=
        \begin{cases}
          1, & |r_{i,\tau}| \le \delta,\\
          \delta/|r_{i,\tau}|, & |r_{i,\tau}|>\delta
        \end{cases}
      \]
      大残差自动降权，增强抗遮挡能力
    \end{block}
    \column{0.48\textwidth}
    \begin{block}{可行域约束}
      \[
        \bm p_t\in\Omega,\quad \forall t
      \]
      栅格地图定义可行域，越界点投影到最近可行栅格
    \end{block}
  \end{columns}
\end{frame}

\subsection{优化推导}

\begin{frame}{梯度与高斯牛顿近似}
  \[
    \frac{\partial r_i}{\partial x}=\frac{\Delta x_i}{d_i},\qquad
    \frac{\partial r_i}{\partial y}=\frac{\Delta y_i}{d_i}
  \]
  \[
    \nabla F(\bm p)=2\sum_i \omega_i r_i
    \begin{bmatrix}
      \Delta x_i/d_i\\
      \Delta y_i/d_i
    \end{bmatrix}
  \]
  \begin{block}{近似海森矩阵}
    \[
      \bm H \approx 2\sum_i \omega_i \bm j_i\bm j_i^{\T},
      \quad \bm j_i=\begin{bmatrix}\Delta x_i/d_i\\ \Delta y_i/d_i\end{bmatrix}
    \]
  \end{block}
\end{frame}

\begin{frame}{滑动窗口变量与收敛条件}
  窗口长度 $w$ 时，优化变量
  \[
    \bm X=[\bm p_{t-w+1}^{\T},\dots,\bm p_t^{\T}]^{\T}\in\R^{2w}
  \]
  \begin{itemize}
    \item 观测项：块对角结构
    \item 运动约束：相邻块耦合，呈稀疏块带状
    \item 收敛判据：$\|\Delta \bm X\|_2/\|\bm X\|_2<10^{-5}$
  \end{itemize}
\end{frame}

\section{算法设计}

\subsection{初始化}

\begin{frame}{初始轨迹构造}
  \begin{enumerate}
    \setlength{\itemsep}{6pt}
    \item \textbf{单帧粗定位}：选信号最强 $K$ 个信标，网格搜索得 $\tilde{\bm p}_t$
    \item \textbf{轨迹修正}：中值滤波去跳点 + 速度上界裁剪异常位移
    \item \textbf{可行域投影}：$\bm p_t^{(0)}=\Pi_{\Omega}(\tilde{\bm p}_t)$
    \item \textbf{缺失补点}：$\bm p_t^{(0)}=\bm p_{t-1}^{(0)}+\alpha(\bm p_{t-1}^{(0)}-\bm p_{t-2}^{(0)})$
  \end{enumerate}
  \vspace{3mm}
  充分初始化可使主优化迭代次数降低约 30\%
\end{frame}

\subsection{求解流程}

\begin{frame}{总体求解流程}
  \begin{enumerate}
    \setlength{\itemsep}{8pt}
    \item 读取当前窗口 RSSI 数据并做异常预处理
    \item 构造初始轨迹 $\bm X^{(0)}$
    \item 执行 IRWLS 迭代求解
    \item 输出当前时刻位置，滑动到下一窗口
  \end{enumerate}
  \vspace{4mm}
  \begin{block}{失败保护机制}
    迭代上限、阻尼缩放、异常回退、信标不足时短时外推
  \end{block}
\end{frame}

\begin{frame}{滑动窗口鲁棒定位算法（IRWLS）}
  \small
  \begin{enumerate}
    \setlength{\itemsep}{4pt}
    \item 输入：窗口观测、初始轨迹、$\lambda_s,\lambda_h,\delta$
    \item \textbf{For} $k=1$ to MaxIter
    \item \quad 计算残差 $r_i$，按 Huber 规则更新 $\omega_i$
    \item \quad 构造法方程 $\bm A\Delta\bm X=\bm b$ 并求解
    \item \quad 更新 $\bm X\leftarrow\bm X+\Delta\bm X$，执行可行域投影
    \item \quad 若 $\|\Delta\bm X\|/\|\bm X\|<\varepsilon$，终止
    \item 输出窗口末端位置估计
  \end{enumerate}
  \normalsize
  复杂度：$\mathcal O(w^3)$；默认 $w=8$，单帧毫秒级
\end{frame}

\begin{frame}{核心参数设置}
  \begin{table}
    \centering
    \caption{推荐参数范围}
    \small
    \begin{tabular}{lll}
      \toprule
      参数 & 建议范围 & 说明 \\
      \midrule
      窗口长度 $w$ & 6--10 & 平衡时延与平滑性 \\
      Huber 阈值 $\delta$ & 2--4 m & 异常抑制强度 \\
      $\lambda_s$ & 0.1--1.0 & 速度平滑权重 \\
      $\lambda_h$ & 0.01--0.2 & 转向平滑权重 \\
      \bottomrule
    \end{tabular}
  \end{table}
  \begin{itemize}
    \item $\lambda_s$ 过小抖动，过大抑制转弯
    \item $\delta$ 过大抑制不足，过小误伤有效观测
  \end{itemize}
\end{frame}

\section{实验结果}

\subsection{主实验}

\begin{frame}{实验设置}
  \begin{columns}[T]
    \column{0.55\textwidth}
    \begin{itemize}
      \setlength{\itemsep}{5pt}
      \item 场景：$120\,\mathrm{m}\times 80\,\mathrm{m}$ 地下停车场
      \item 信标：16 个固定 BLE 信标
      \item 轨迹：直行、转弯、绕障，共 300 s，1 Hz 采样
      \item 噪声：高斯扰动 + 随机遮挡
    \end{itemize}
    \column{0.42\textwidth}
    \begin{block}{对比方法}
      \begin{enumerate}
        \item RSSI-Only
        \item KF-Baseline
        \item Proposed（本文）
      \end{enumerate}
    \end{block}
  \end{columns}
\end{frame}

\begin{frame}{主实验结果对比}
  \begin{table}
    \centering
    \caption{定位精度与实时性对比}
    \begin{tabular}{lcccc}
      \toprule
      方法 & MAE/m & RMSE/m & 方向方差 & 耗时/ms \\
      \midrule
      RSSI-Only & 4.83 & 6.21 & 0.198 & 0.6 \\
      KF-Baseline & 3.15 & 4.02 & 0.093 & 1.1 \\
      \textbf{Proposed} & \textbf{2.08} & \textbf{2.76} & \textbf{0.041} & 2.4 \\
      \bottomrule
    \end{tabular}
  \end{table}
  \begin{itemize}
    \item 融合模型精度与轨迹连续性明显优于基线
    \item 单帧计算 2.4 ms，满足应急实时需求
  \end{itemize}
\end{frame}

\begin{frame}{遮挡鲁棒性与窗口长度}
  \begin{columns}[T]
    \column{0.48\textwidth}
    \begin{block}{遮挡实验}
      \begin{itemize}
        \item 遮挡比例 0\% $\rightarrow$ 40\%
        \item RSSI-Only 误差快速上升
        \item 30\% 遮挡下本文方法仍约 3 m
      \end{itemize}
    \end{block}
    \column{0.48\textwidth}
    \begin{block}{窗口长度}
      \begin{itemize}
        \item $w$ 从 4 增至 12：误差先降后稳
        \item 耗时线性增长
        \item 默认取 $w=8$
      \end{itemize}
    \end{block}
  \end{columns}
\end{frame}

\subsection{扩展分析}

\begin{frame}{灵敏度分析}
  \begin{table}
    \centering
    \small
    \begin{tabular}{lcc}
      \toprule
      因素变化 & RMSE 变化率 & 结论 \\
      \midrule
      噪声翻倍 & +42\% & 鲁棒算法可控 \\
      信标减半 & +67\% & 几何覆盖是关键 \\
      $n$ 偏差 0.3 & +38\% & 标定精度重要 \\
      \bottomrule
    \end{tabular}
  \end{table}
  \begin{itemize}
    \item 8$\rightarrow$12 个信标阶段误差下降明显，超过 16 个收益递减
    \item 车道转角区域局部加密优于全局均匀加密
  \end{itemize}
\end{frame}

\begin{frame}{不确定性分析}
  \begin{block}{局部协方差近似}
    \[
      \bm\Sigma_X\approx (\bm J^\T\bm W\bm J)^{-1}
    \]
  \end{block}
  \begin{block}{二维置信椭圆}
    \[
      (\bm p-\hat{\bm p})^\T \bm\Sigma_t^{-1}(\bm p-\hat{\bm p})\le \chi^2_{2,p}
    \]
  \end{block}
  \begin{itemize}
    \item 误差来源：RSSI 噪声、路径损耗标定误差、几何遮挡偏差
    \item 在线输出：椭圆轴长、有效信标数、残差统计、异常保护状态
  \end{itemize}
\end{frame}

\begin{frame}{案例分析与应用价值}
  \begin{enumerate}
    \setlength{\itemsep}{6pt}
    \item \textbf{火灾疏散引导}：实时位置 + 可信度，判断偏离疏散通道
    \item \textbf{救援协同定位}：多终端独立定位，跳变时可信度告警
    \item \textbf{信标局部失效}：30\% 失效时运动约束维持连续轨迹
  \end{enumerate}
  \vspace{4mm}
  调度支持：路径推荐、位置回放、覆盖评估、风险预警
\end{frame}

\section{总结}

\begin{frame}{主要工作与结论}
  \begin{columns}[T]
    \column{0.48\textwidth}
    \begin{block}{主要工作}
      \begin{itemize}
        \setlength{\itemsep}{4pt}
        \item BLE RSSI + 步行约束融合模型
        \item 鲁棒 IRWLS + 滑动窗口优化
        \item 初始化与可行域投影机制
        \item 多噪声、多遮挡实验验证
      \end{itemize}
    \end{block}
    \column{0.48\textwidth}
    \begin{block}{主要结论}
      \begin{itemize}
        \setlength{\itemsep}{4pt}
        \item 显著优于单一 RSSI 定位
        \item 运动约束抑制轨迹抖动
        \item 遮挡场景表现稳定
        \item 满足实时计算需求
      \end{itemize}
    \end{block}
  \end{columns}
\end{frame}

\begin{frame}{不足与后续工作}
  \begin{block}{不足}
    \begin{itemize}
      \item 尚未在真实地下停车场长周期实测
      \item 多目标互扰场景考虑不足
      \item 多层立体拓扑尚未引入
    \end{itemize}
  \end{block}
  \begin{block}{后续方向}
    真实演练数据采集；融合 IMU/UWB 因子图；多目标协同定位；可视化指挥端联调
  \end{block}
\end{frame}

\begin{frame}{参考文献}
  \small
  \begin{enumerate}
    \setlength{\itemsep}{3pt}
    \item J. Nocedal, S. J. Wright, \emph{Numerical Optimization}, Springer, 2006.
    \item R. Faragher, R. Harle, BLE beacons fingerprinting, \emph{IEEE JSAC}, 2015.
    \item B. Li \emph{et al.}, Robust RSSI-based indoor localization, \emph{Sensors}, 2021.
    \item S. Boyd, L. Vandenberghe, \emph{Convex Optimization}, Cambridge, 2004.
    \item 陈务深, 室内定位与优化计算课程讲义, 2026.
  \end{enumerate}
\end{frame}

\begin{frame}
  \begin{center}
    {\Huge\calligra Thanks for your attention!}
    \vspace{1cm}
    {\Huge Q \& A}
  \end{center}
\end{frame}

\end{document}
